% ============================================================================
%  The gravitational N-body problem — report (English)
%  Build: python tools/build_documents.py   (needs xelatex)
% ============================================================================
\documentclass[11pt,a4paper]{article}

\usepackage{amsmath,amssymb}
\usepackage{fontspec}
\usepackage{unicode-math}
\setmathfont{Latin Modern Math}
\usepackage{polyglossia}
\setdefaultlanguage{english}
\usepackage[margin=2.3cm]{geometry}
\usepackage{graphicx}
\usepackage{booktabs}
\usepackage{siunitx}
\sisetup{exponent-product=\times, group-minimum-digits=5, group-digits=integer}
\usepackage[font=small,labelfont=bf,margin=0.5cm]{caption}
\usepackage{float}
\usepackage{xcolor}
\definecolor{link}{HTML}{1f5fae}
\usepackage[colorlinks=true,linkcolor=link,citecolor=link,urlcolor=link]{hyperref}
\hypersetup{pdftitle={Jesús García Izquierdo - https://github.com/jesus-garcia-izquierdo/N-body-problem}, pdfauthor={Jesús García Izquierdo}}
\usepackage{enumitem}
\setlist{nosep,leftmargin=1.4em}
\usepackage{fancyhdr}
\pagestyle{fancy}
\fancyhf{}
\fancyhead[L]{\small The gravitational $N$-body problem}
\fancyhead[R]{\small J. García Izquierdo}
\fancyfoot[C]{\small\thepage}
\renewcommand{\headrulewidth}{0.3pt}
\setlength{\parskip}{0.35em}

\graphicspath{{figures/}}
\input{numbers.tex}

\newcommand{\vect}[1]{\mathbf{#1}}
\newcommand{\dd}{\mathrm{d}}
\newcommand{\repo}{\href{https://github.com/jesus-garcia-izquierdo/N-body-problem}{github.com/jesus-garcia-izquierdo/N-body-problem}}

\begin{document}

% ---------------------------------------------------------------------------
\begin{titlepage}
\centering
\vspace*{2.5cm}
{\Large\bfseries Jesús García Izquierdo - {\urlstyle{same}\hypersetup{urlcolor=black}\url{https://github.com/jesus-garcia-izquierdo/N-body-problem}}\par}
\vspace{0.6cm}
{\LARGE\bfseries The gravitational $N$-body problem\par}
\vspace{0.25cm}
{\large Symplectic integration, conservation laws, stability and chaos\par}
\vspace{1.5cm}
{\normalsize October 2026\par}
\vspace{1.2cm}
{\small Python code, videos and data: \repo\par}
\vspace{2cm}
\begin{minipage}{0.86\textwidth}
\small
\textbf{Abstract.} We study the motion of $N$ point masses interacting only through
Newtonian gravity. After reviewing the physics (equations of motion, Hamiltonian
structure, conservation laws, the virial theorem, linear stability and Lyapunov
exponents) we integrate the equations numerically with the velocity Verlet method and
four alternative schemes, using a fully vectorized Python implementation. Eight
experiments are analysed: the Sun--Earth--Moon system, a four-body system with a
massive central body, a chaotic six-body system, a comparison of integrators, the
figure-eight choreography, the stability of Lagrange's equilateral triangle, the Solar
System and the gravitational collapse of a star cluster. Energy is conserved to
\SemErrE\ (Sun--Earth--Moon, two years), linear and angular momentum to machine
precision, the measured orders of convergence match theory, the instability rate of
Lagrange's triangle agrees with linear theory to three digits
(\LagSigma\ vs.\ \LagSigmaTh), Kepler's third law is recovered with slope \SolarSlope,
and the six-body system has a Lyapunov time of only \ChaosLyapDays\ days.
\end{minipage}
\vfill
{\footnotesize This report extends and corrects the 2023 project
\emph{Problema de los n-cuerpos} by G.~Haro Pérez and J.~García Izquierdo
(MATLAB implementation).\par}
\end{titlepage}

\tableofcontents
\newpage

% ===========================================================================
\section{Introduction}
% ===========================================================================

The $N$-body problem asks for the motion of $N$ masses that attract each other
gravitationally. For $N=2$ it was solved by Newton: the bodies move on conic sections.
For $N\ge 3$ no general closed-form solution exists. Bruns (1887) showed that there are
no further algebraic integrals beyond the ten classical ones, and Poincaré (1890)
showed that the motion can be extremely sensitive to the initial conditions, the
phenomenon we now call \emph{chaos}. Sundman (1912) did find a convergent power series
for $N=3$, but it converges far too slowly to be of any practical use.

The only general way forward is therefore numerical integration. A numerical solution,
however, is only trustworthy if it respects what the exact dynamics guarantees:
conservation of energy, momentum and angular momentum, and statistical relations such as
the virial theorem. This report has three goals:
\begin{enumerate}
  \item to present the physics needed to interpret a simulation (Section~\ref{sec:model});
  \item to explain why a \emph{symplectic} integrator such as velocity Verlet is the
        right tool, and to quantify its accuracy (Section~\ref{sec:numerics});
  \item to analyse eight experiments that range from the Solar System to chaos and
        star-cluster formation (Section~\ref{sec:results}).
\end{enumerate}
Every number quoted in the text is produced by the code in the repository; the script
\texttt{tools/build\_documents.py} regenerates all figures and values.

% ===========================================================================
\section{Physical model}\label{sec:model}
% ===========================================================================

\subsection{Equations of motion}

Let body $i$ have mass $m_i$, position $\vect r_i(t)$ and velocity $\vect v_i(t)$.
Newton's law of gravitation together with the superposition principle gives
\begin{equation}
  \ddot{\vect r}_i = G\sum_{j\ne i} m_j\,\frac{\vect r_j-\vect r_i}{|\vect r_j-\vect r_i|^{3}},
  \qquad i=1,\dots,N,
  \label{eq:eom}
\end{equation}
with $G=\SI{6.67430e-11}{m^3.kg^{-1}.s^{-2}}$. These are $3N$ coupled second-order
ordinary differential equations, or $6N$ first-order ones for $(\vect r_i,\vect v_i)$.
Given positions and velocities at $t=0$ the solution exists and is unique as long as no
two bodies collide.

\subsection{Hamiltonian structure}

Writing the energy in terms of positions $\vect q_i=\vect r_i$ and momenta
$\vect p_i=m_i\vect v_i$,
\begin{equation}
  H(\vect q,\vect p)=\underbrace{\sum_i\frac{|\vect p_i|^2}{2m_i}}_{T(\vect p)}
  \;\underbrace{-\sum_{i<j}\frac{G m_i m_j}{|\vect q_i-\vect q_j|}}_{V(\vect q)},
  \label{eq:hamiltonian}
\end{equation}
Hamilton's equations $\dot{\vect q}_i=\partial H/\partial\vect p_i$,
$\dot{\vect p}_i=-\partial H/\partial\vect q_i$ reproduce Eq.~\eqref{eq:eom}. Two
properties of the exact flow will guide the choice of integrator: it preserves volume
(indeed the symplectic two-form $\sum_i \dd\vect q_i\wedge \dd\vect p_i$) in the
$6N$-dimensional phase space (Liouville's theorem), and it is time-reversible. Note also
that $H=T(\vect p)+V(\vect q)$ is \emph{separable}, which is what makes the Verlet
splitting possible (Section~\ref{sec:verlet}).

\subsection{Conservation laws}

By Noether's theorem, the symmetries of an isolated system imply conserved quantities:
\begin{align}
  \text{translations} &\;\Rightarrow\; \vect P=\sum_i m_i\vect v_i, \\
  \text{rotations}    &\;\Rightarrow\; \vect L=\sum_i \vect r_i\times m_i\vect v_i, \\
  \text{time shifts}  &\;\Rightarrow\; E=K+U,\quad K=\tfrac12\sum_i m_i v_i^2,\quad
                        U=-\sum_{i<j}\frac{G m_i m_j}{r_{ij}} .
  \label{eq:energy}
\end{align}
The proofs show what a numerical method must respect. $\dot{\vect P}=0$ follows from
$\vect F_{ij}=-\vect F_{ji}$ (Newton's third law), so any scheme built on exactly
antisymmetric pair forces conserves $\vect P$ up to round-off. $\dot{\vect L}=0$
follows because the forces are central, $\vect F_{ij}\parallel\vect r_i-\vect r_j$.
The potential energy counts every pair once; summing over all ordered pairs and dividing
by two, as in the original report, is equivalent.

Since $\vect P$ is constant, the centre of mass (CM) moves uniformly with
$\vect V_\text{cm}=\vect P/M$. We work in the CM frame ($\vect P=0$) whenever possible:
the system does not drift away, and the kinetic energy loses the trivial term
$\tfrac12 M V_\text{cm}^2$, which is essential for the virial theorem.

\subsection{The virial theorem}\label{sec:virial}

Define the virial $\mathcal G=\sum_i\vect p_i\cdot\vect r_i$. Its derivative is
$\dot{\mathcal G}=2K+\sum_i\vect F_i\cdot\vect r_i$. Because $U$ is homogeneous of degree
$-1$ in the positions, Euler's theorem gives $\sum_i\vect r_i\cdot\nabla_i U=-U$, i.e.\
$\sum_i\vect F_i\cdot\vect r_i=U$. Averaging over a time $\tau$,
\begin{equation}
  \frac{\mathcal G(\tau)-\mathcal G(0)}{\tau}=2\langle K\rangle+\langle U\rangle
  \;\xrightarrow[\tau\to\infty]{}\;0
  \quad\Longrightarrow\quad
  Q\equiv\frac{2\langle K\rangle}{-\langle U\rangle}=1,
  \label{eq:virial}
\end{equation}
provided the motion is bounded. Three conditions are therefore required: a bound system
($E<0$) with no escapers, $K$ measured in the CM frame, and a long averaging time. In
equilibrium $E=\tfrac12\langle U\rangle=-\langle K\rangle$.

\subsection{The two-body problem}

With $\vect R$ the CM and $\vect r=\vect r_2-\vect r_1$, the two-body problem reduces to
one body of reduced mass $\mu=m_1m_2/(m_1+m_2)$ in a central field. The orbit is a conic;
for bound orbits two results are used repeatedly below:
\begin{equation}
  v^2=GM\left(\frac{2}{r}-\frac1a\right)\ \text{(vis-viva)},\qquad
  T^2=\frac{4\pi^2}{GM}\,a^3\ \text{(Kepler's third law)},
\end{equation}
where $a$ is the semi-major axis. A circular orbit has $v_c=\sqrt{GM/r}$.

\subsection{Stability, chaos and Lyapunov exponents}\label{sec:chaos}

Let $\boldsymbol\delta(t)$ be the phase-space separation between two solutions that start
a distance $\delta_0$ apart. For almost every initial direction
\begin{equation}
  \lambda_\text{max}=\lim_{t\to\infty}\lim_{\delta_0\to0}\frac1t\ln\frac{|\boldsymbol\delta(t)|}{\delta_0}
\end{equation}
exists and is the \emph{maximal Lyapunov exponent}. Regular motion has
$\lambda_\text{max}=0$ (separations grow at most polynomially); chaotic motion has
$\lambda_\text{max}>0$, and errors grow by a factor $e$ every Lyapunov time
$t_L=1/\lambda_\text{max}$. With an initial uncertainty $\delta_0$ and a tolerance $\Delta$,
predictions are possible only up to $t_\text{pred}\approx t_L\ln(\Delta/\delta_0)$: improving
the measurements a thousandfold buys only $\ln 1000\approx7$ extra Lyapunov times.

We estimate $\lambda_\text{max}$ with Benettin's method: a ``shadow'' trajectory is
integrated alongside the reference one; every $\tau$ its separation $d_k$ is measured,
$\ln(d_k/\delta_0)$ is accumulated and the shadow is rescaled back to distance $\delta_0$
along the same direction, so that $\lambda\approx t^{-1}\sum_k\ln(d_k/\delta_0)$.
Positions and velocities are made dimensionless with the rms radius $L$ and the time
$L/v_\text{rms}$.

\emph{Linear stability} is the local version of the same idea: a solution is linearly
unstable if small perturbations grow exponentially. Lagrange's equilateral solution
(Section~\ref{sec:lagrange}) is a textbook example where the growth rate can be computed
analytically and compared with the simulation.

\subsection{Softening and units}

For star clusters, where particles represent many stars and close two-body encounters
are not of interest, we replace $1/r\to1/\sqrt{r^2+\varepsilon^2}$ (Plummer softening).
The force remains the gradient of the softened potential, so the softened energy is
still conserved. Abstract problems use $N$-body units $G=M=1$, in which all quantities
are of order one; physical scales are recovered by noting that time scales as
$\sqrt{L^3/GM}$.

% ===========================================================================
\section{Numerical method}\label{sec:numerics}
% ===========================================================================

\subsection{Velocity Verlet}\label{sec:verlet}

Writing $\vect a_n=\vect a(\vect r_n)$ for the accelerations of Eq.~\eqref{eq:eom}, one
step of length $h$ is
\begin{align}
  \vect r_{n+1}&=\vect r_n+h\,\vect v_n+\tfrac{h^2}{2}\,\vect a_n, \label{eq:verlet1}\\
  \vect a_{n+1}&=\vect a(\vect r_{n+1}), \\
  \vect v_{n+1}&=\vect v_n+\tfrac{h}{2}\left(\vect a_n+\vect a_{n+1}\right). \label{eq:verlet3}
\end{align}
The position update is a second-order Taylor expansion and the velocity update is the
trapezoidal rule for $\int\vect a\,\dd t$. Introducing $\vect v_{n+1/2}=\vect v_n+\tfrac h2\vect a_n$
the same algorithm reads ``half kick -- drift -- half kick''. This is the operator
splitting $e^{\frac h2 V}e^{hT}e^{\frac h2 V}$ of the Hamiltonian~\eqref{eq:hamiltonian}:
the drift is the exact flow of $T$ and the kicks are the exact flow of $V$. Hence
\begin{itemize}
  \item the method is \textbf{symplectic} (a composition of Hamiltonian flows) and
        \textbf{time-reversible} (the composition is symmetric);
  \item by the Baker--Campbell--Hausdorff formula the local error is $\mathcal O(h^3)$ and
        the global error $\mathcal O(h^2)$;
  \item it conserves $\vect P$ and $\vect L$ exactly (each sub-step does);
  \item it needs one force evaluation per step, because $\vect a_{n+1}$ is reused as the
        $\vect a_n$ of the next step (the original MATLAB code evaluated it twice).
\end{itemize}

\textbf{Why symplectic matters.} Backward error analysis shows that a symplectic method
of order $p$ follows \emph{exactly} the dynamics of a modified, ``shadow'' Hamiltonian
$\tilde H=H+h^pH_p+\dots$, up to exponentially small terms. The numerical trajectory
conserves $\tilde H$, so the error in $H$ stays bounded by $\mathcal O(h^p)$ over
exponentially long times instead of drifting. This is the single most important property
for long orbital integrations.

\subsection{Other integrators}

To show the effect of order and structure, four other schemes are implemented
(Table~\ref{tab:integrators}). Explicit Euler expands phase-space areas by
$1+\omega^2h^2$ per step on a harmonic oscillator, so the energy grows exponentially;
swapping two lines (symplectic Euler) makes the area factor exactly one. Yoshida's method
composes three Verlet steps of lengths $w_1h,\,w_0h,\,w_1h$ with
$w_1=1/(2-2^{1/3})$ and $w_0=-2^{1/3}/(2-2^{1/3})$ (note the step backwards in time),
which cancels the $h^3$ error term. RK4 is fourth order but not symplectic.

\begin{table}[H]
\centering
\caption{Integrators implemented in \texttt{nbody/integrators.py}.}
\label{tab:integrators}
\begin{tabular}{lcccc}
\toprule
Method & Order & Symplectic & Reversible & Forces per step\\
\midrule
Explicit Euler    & 1 & no  & no  & 1\\
Symplectic Euler  & 1 & yes & no  & 1\\
Velocity Verlet   & 2 & yes & yes & 1\\
Yoshida-4         & 4 & yes & yes & 3\\
Runge--Kutta 4    & 4 & no  & no  & 4\\
\bottomrule
\end{tabular}
\end{table}

\subsection{Implementation}

The code is a small Python package (\texttt{nbody}) built on NumPy. The state is stored
as arrays of shape $(N,3)$ (body $\times$ component) and a trajectory as arrays of shape
$(T,N,3)$ (time $\times$ body $\times$ component). Accelerations are computed without any
Python loop over bodies: broadcasting builds the $(N,N,3)$ tensor of separations
$\vect r_j-\vect r_i$, the diagonal is excluded by setting $r_{ii}^2=\infty$, and an
Einstein summation contracts the index $j$. The cost is $\mathcal O(N^2)$ per step, as
for any direct-summation code. Memory for the trajectory is allocated once (the MATLAB
version grew its arrays at every step, an $\mathcal O(T^2)$ cost). Typical run times on a
laptop are a second for the Sun--Earth--Moon system ($6.3\times10^4$ steps) and a few
seconds for the Solar System ($1.2\times10^5$ steps).

\subsection{Diagnostics and choice of time step}

Conservation is measured \emph{relative to the initial value}:
$|E(t)-E_0|/|E_0|$ for the energy, and for the vectors $\vect P$ and $\vect L$ the norm
of the change divided by the natural scales $\max_t\sum_i m_i|\vect v_i|$ and
$\max_t\sum_i m_i|\vect r_i||\vect v_i|$ (needed because $\vect P=0$ in the CM frame).
The time step is set by the shortest dynamical time of the system, with the rule of
thumb $h\lesssim t_\text{dyn}/100$, and every result was checked by halving $h$.

% ===========================================================================
\section{Results}\label{sec:results}
% ===========================================================================

\subsection{The Sun--Earth--Moon system}

The bodies start in the ecliptic plane with the Earth at aphelion and the Moon at apogee
(Table~\ref{tab:sem}). The shortest period is the Moon's, 27.3 days, so $h=\SI{1000}{s}$
gives about 2400 steps per lunar orbit.

\begin{table}[H]
\centering
\caption{Initial conditions of the Sun--Earth--Moon system (before moving to the CM frame).
Integration time 2 years, $h=\SI{1000}{s}$.}
\label{tab:sem}
\begin{tabular}{lccc}
\toprule
 & Sun & Earth & Moon\\
\midrule
Mass (kg) & \num{1.98847e30} & \num{5.9722e24} & \num{7.342e22}\\
$x_0$ (m) & 0 & \num{1.52098232e11} (aphelion) & $x_\oplus+\num{4.055e8}$ (apogee)\\
$v_{y,0}$ (m/s) & 0 & \num{29290} & $v_\oplus+\num{964}$\\
\bottomrule
\end{tabular}
\end{table}

Figure~\ref{fig:sem-orbits} shows the orbits. At heliocentric scale the Moon's orbit is
invisible (the Earth--Moon distance is 1/375 of the Sun--Earth distance), so it is also
shown relative to the Earth: the lunar ellipse visibly \emph{precesses}, a direct effect
of the Sun's tidal perturbation that a two-body model would miss.

\begin{figure}[H]
\centering
\includegraphics[width=\textwidth]{sem_orbits}
\caption{Sun--Earth--Moon system over two years. Left: orbits about the centre of mass
(the Earth and the Moon overlap at this scale). Right: the Moon relative to the Earth; its
orbit precesses under the Sun's perturbation.}
\label{fig:sem-orbits}
\end{figure}

The conservation laws are shown in Fig.~\ref{fig:sem-cons}. The energy error oscillates
with the orbital period and stays below \SemErrE: there is no secular drift, as expected
from a symplectic method. The linear and angular momenta are conserved to
\SemErrP\ and \SemErrL\ respectively, i.e.\ to machine precision; the slow growth is the
accumulation of round-off errors. The time averages give
$\langle K\rangle=\SemKmean\,\mathrm{J}$ and $-\tfrac12\langle U\rangle=\SemUhalf\,\mathrm{J}$, so
$Q=\SemVirial$: the virial theorem holds to better than $10^{-4}$.

\begin{figure}[H]
\centering
\includegraphics[width=\textwidth]{sem_conservation}
\caption{Conservation laws for the Sun--Earth--Moon run: (a) kinetic, potential and total
energy; (b) relative energy error, bounded and oscillating; (c) relative errors of
$\vect P$ and $\vect L$, at the level of round-off.}
\label{fig:sem-cons}
\end{figure}

If the Sun is kept at rest initially (as in the original report) the total momentum is
not zero and the CM moves at $V_\text{cm}=\SemVcm$~m/s. Here this is negligible because
the Sun dominates the mass, but in general it biases $\langle K\rangle$ and therefore $Q$.

\subsection{Four bodies around a massive central body}

This system reproduces Section~4.2 of the original report (Table~\ref{tab:four}): a
central body of \SI{2e40}{kg} (about $10^{10}$ solar masses, comparable to a galactic
nucleus), two stars of about one solar mass in the $xy$ plane, and a light body on a
polar orbit.

\begin{table}[H]
\centering
\caption{Initial conditions of the four-body system ($d=\SI{1.496e14}{m}$,
$M_\odot=\SI{1.9891e30}{kg}$). Integration time 15 years, $h=\SI{1000}{s}$.}
\label{tab:four}
\begin{tabular}{lcccc}
\toprule
 & Body 1 & Body 2 & Body 3 & Body 4\\
\midrule
Mass (kg) & $2.02\,M_\odot$ & $0.98\,M_\odot$ & \num{2e40} & \num{2e20}\\
$\vect r_0$ (m) & $(0,0,0)$ & $(d,0,0)$ & $(d/2,0,0)$ & $(d/2,0,10^{14})$\\
$\vect v_0$ & $(0,-\sqrt{Gm_3/d},0)$ & $(0,\sqrt{Gm_3/d},0)$ & $\vect 0$ & $(-\sqrt{Gm_3/10^{14}},0,0)$\\
\bottomrule
\end{tabular}
\end{table}

The original report intended circular orbits, but the speeds use the distance $d$, which is
\emph{twice} the actual distance $r=d/2$ to the central body. Then $v^2=GM/(2r)$ and the
vis-viva equation gives $a=2r/3$: the orbits are ellipses with eccentricity $e=0.5$ whose
apocentre is the starting point. The simulation measures $e=\FourEcc$, confirming the
analysis. Body 4 starts at the correct distance and moves on a circle.

\begin{figure}[H]
\centering
\includegraphics[width=0.82\textwidth]{four_orbits}
\caption{Orbits of the four-body system over 15 years (about 250 orbits of bodies 1 and 2).
The ellipses are closed and stable; the massive body barely moves.}
\label{fig:four-orbits}
\end{figure}

\begin{figure}[H]
\centering
\includegraphics[width=\textwidth]{four_conservation}
\caption{Conservation laws for the four-body system. The energy error peaks at pericentre
passages, where the time step is relatively largest, but does not grow.}
\label{fig:four-cons}
\end{figure}

With eccentric orbits the energy error peaks at each pericentre passage, where the orbital
time scale is shortest, and reaches \FourErrE\ (Fig.~\ref{fig:four-cons}); yet the
envelope is flat over 15 years. $\vect P$ and $\vect L$ are conserved to \FourErrP\ and
\FourErrL, and the virial ratio is $Q=\FourVirial$.

\subsection{Chaos in a six-body system}

Six bodies of \SI{1e10}{kg} start on the $\pm x$, $\pm y$, $\pm z$ semi-axes at
\SI{1e5}{m}, each with the speed ($\approx\SI{2.1}{mm/s}$) of a circular orbit around the
next one (Table~\ref{tab:six}). We integrate four versions for 10 years: the reference and
three runs in which a single quantity is changed by 0.1\,\%.

\begin{table}[H]
\centering
\caption{Six-body system ($u=\SI{1.5e-3}{m/s}$). Integration time 10 years, $h=\SI{1000}{s}$.}
\label{tab:six}
\begin{tabular}{lcccccc}
\toprule
 & 1 & 2 & 3 & 4 & 5 & 6\\
\midrule
$\vect r_0/\SI{e5}{m}$ & $(1,0,0)$ & $(0,1,0)$ & $(0,0,1)$ & $(-1,0,0)$ & $(0,-1,0)$ & $(0,0,-1)$\\
$\vect v_0/u$ & $(1,1,0)$ & $(0,1,1)$ & $(-1,0,1)$ & $(-1,-1,0)$ & $(0,-1,-1)$ & $(1,0,-1)$\\
\bottomrule
\end{tabular}
\end{table}

\begin{figure}[H]
\centering
\includegraphics[width=\textwidth]{chaos_orbits}
\caption{Six-body system: reference run and three runs with a 0.1\,\% change in $x_1$,
$v_{y,2}$ and $m_3$. Bodies pair up into binaries and some are ejected, differently in
each run.}
\label{fig:chaos-orbits}
\end{figure}

Figure~\ref{fig:chaos-orbits} shows that the final configurations are qualitatively
different. Figure~\ref{fig:chaos-div} quantifies this: the separation from the reference
grows exponentially (a straight line on a log scale) after the first close encounter,
around $t\approx2$ years, until it saturates at the size of the system.

The most instructive curve is the last one: the \emph{same} initial condition integrated
with $h$ and with $h/2$ also diverges exponentially, reaching \ChaosHalfFinal~m after 10
years. The truncation error acts as one more perturbation, so beyond a few Lyapunov times
neither run reproduces the exact trajectory. This does not invalidate the simulation:
$E$, $\vect P$ and $\vect L$ are still conserved (\ChaosErrE, \ChaosErrP, \ChaosErrL), and
the shadowing lemma guarantees that the numerical orbit stays close to \emph{some} exact
orbit with slightly different initial conditions. Statistical properties are reliable;
individual trajectories are not.

\begin{figure}[H]
\centering
\includegraphics[width=\textwidth]{chaos_divergence}
\caption{Separation $\delta r(t)$ between each perturbed run and the reference. The
yellow curve compares the same initial condition integrated with $h$ and $h/2$.}
\label{fig:chaos-div}
\end{figure}

Benettin's method (Fig.~\ref{fig:chaos-lyap}) gives
$\lambda_\text{max}\approx\ChaosLambdaYr~\text{yr}^{-1}$ after 10 years, i.e.\ a Lyapunov
time of about \ChaosLyapDays\ days. The estimate jumps after the first close encounter and
is still drifting slowly, a sign that the chaotic phase starts at $t\approx2$~yr. Over
ten years an initial error is amplified by roughly \ChaosAmpExp\ orders of magnitude: not
even double-precision initial conditions (relative error $10^{-16}$) would allow a
deterministic prediction.

\begin{figure}[H]
\centering
\includegraphics[width=\textwidth]{chaos_lyapunov}
\caption{Running estimate of the maximal Lyapunov exponent (Benettin's method,
renormalization every 20 steps, $\delta_0=10^{-8}$). The dashed line is the average over
the last 20\,\% of the run.}
\label{fig:chaos-lyap}
\end{figure}

\subsection{Comparison of integrators}

To isolate the effect of the integrator we use the figure-eight orbit of
Section~\ref{sec:eight}, with $G=m=1$ and period $T\approx6.326$. Figure~\ref{fig:int-energy}
shows the maximum energy error per period over 200 periods with a deliberately coarse
step, $h=T/300$. Explicit Euler gains energy at every step and the system disintegrates
within ten periods. RK4 starts with a tiny error but drifts \emph{linearly}. The three
symplectic schemes show no drift at all, as the shadow-Hamiltonian argument predicts; at
equal cost Verlet is far more robust than its formal order suggests, and Yoshida-4
combines high order with bounded error.

\begin{figure}[H]
\centering
\includegraphics[width=\textwidth]{integrators_energy}
\caption{Long-term energy error (maximum per period) for $h=T/300$ on the figure-eight
orbit. Only the non-symplectic methods show a secular drift.}
\label{fig:int-energy}
\end{figure}

Figure~\ref{fig:int-conv} plots the position error after one period against the number
of force evaluations, the real cost. The fitted slopes give orders
\OrderEuler, \OrderSympEuler, \OrderVerlet, \OrderYoshida\ and \OrderRK, matching the
theoretical values 1, 1, 2, 4 and 4.

\begin{figure}[H]
\centering
\includegraphics[width=\textwidth]{integrators_convergence}
\caption{Order of convergence: position error after one period against the number of
force evaluations. Slopes are fitted on the three smallest steps.}
\label{fig:int-conv}
\end{figure}

\subsection{The figure-eight choreography}\label{sec:eight}

Moore (1993) found numerically, and Chenciner and Montgomery (2000) proved, that three
equal masses can chase each other along a single figure-eight curve, a third of a period
apart. The initial conditions are
$\vect r_1=-\vect r_2=(-0.97000436,\,0.24308753)$, $\vect r_3=\vect 0$,
$\vect v_3=(-0.93240737,\,-0.86473146)$, $\vect v_1=\vect v_2=-\vect v_3/2$.
The total angular momentum is exactly zero, even though the bodies rotate: the two lobes
are travelled in opposite senses.

With $h=T/6000$ (Fig.~\ref{fig:eight}) the bodies return to their initial positions within
\EightRetOne\ after one period and \EightRetTen\ after ten (the initial data themselves are
given to eight digits). The energy error is bounded at \EightErrE\ and $Q=\EightVirial$.

\begin{figure}[H]
\centering
\includegraphics[width=\textwidth]{eight_orbit}
\caption{Figure-eight choreography over ten periods: the three orbits coincide. Right:
bounded relative energy error.}
\label{fig:eight}
\end{figure}

Changing one velocity component by 0.1\,\% produces a separation that grows only slowly
and \emph{without} the exponential trend of the six-body system (Fig.~\ref{fig:eight-stab}),
reaching \EightDivFinal\ after ten periods. The orbit is surrounded by invariant KAM tori:
it is stable.

\begin{figure}[H]
\centering
\includegraphics[width=\textwidth]{eight_stability}
\caption{Separation between the figure-eight orbit and a copy with $v_{x,1}$ changed by
0.1\,\%: slow, non-exponential growth.}
\label{fig:eight-stab}
\end{figure}

\subsection{Stability of Lagrange's triangle}\label{sec:lagrange}

Lagrange (1772) showed that three bodies at the vertices of an equilateral triangle of
side $d$, rotating rigidly with $\omega^2=GM/d^3$, form an exact solution for any masses.
Linearizing in the rotating frame, the perturbation exponents satisfy
\begin{equation}
  \lambda^4+\omega^2\lambda^2+\tfrac{27}{4}\,\beta\,\omega^4=0,\qquad
  \beta=\frac{m_1m_2+m_2m_3+m_3m_1}{M^2}.
\end{equation}
The roots are purely imaginary (stable oscillations) only if $27\beta<1$, which is
Routh's criterion. For the Sun, Jupiter and a Trojan asteroid ($\beta\approx10^{-3}$) the
configuration is stable, which is why Trojan asteroids cluster at Jupiter's L4 and L5
points. For equal masses $\beta=1/3$, the roots satisfy $\lambda^2=\omega^2(-1\pm i\sqrt8)/2$
and the growth rate is $\operatorname{Re}\lambda=\omega/\sqrt2$.

In the simulation no perturbation is added by hand: the truncation error of Verlet is
enough (Fig.~\ref{fig:lagrange}). For equal masses the deformation of the triangle grows
exponentially from $10^{-6}$ to order one, and a linear fit of $\ln\delta$ gives
$\sigma=\LagSigma$, against the theoretical $\LagSigmaTh$. For the Sun--Jupiter--Trojan
masses the deformation stays below \LagTrojanMax. The small periodic deformation at the
start is not an instability: Verlet solves exactly a shadow Hamiltonian whose circular
orbit has a slightly different radius.

\begin{figure}[H]
\centering
\includegraphics[width=\textwidth]{lagrange}
\caption{Deformation of Lagrange's triangle (largest relative change of a side). Equal
masses: exponential growth at the rate predicted by linear theory. Sun--Jupiter--Trojan:
bounded.}
\label{fig:lagrange}
\end{figure}

\subsection{The Solar System and Kepler's third law}

The Sun and the eight planets are initialized with their real masses, semi-major axes,
eccentricities and inclinations (Table~\ref{tab:planets}). Each planet starts at
perihelion with speed $v_p=\sqrt{GM_\odot(1+e)/(a(1-e))}$ and a random orbital phase; its
orbital plane is tilted about the Sun--perihelion line. These are illustrative initial
conditions, not the ephemeris of a given date. The step $h=0.5$~days is set by Mercury
(176 steps per orbit), and the run covers one orbit of Neptune, 165 years.

\begin{table}[H]
\centering
\caption{Planetary parameters and orbital periods measured in the simulation (counting
revolutions of the polar angle about the Sun).}
\label{tab:planets}
\begin{tabular}{lccccc}
\toprule
Planet & Mass (kg) & $a$ (AU) & $e$ & $i$ (°) & $T$ measured (yr)\\
\midrule
Mercury & \num{3.3011e23} & 0.387 & 0.206 & 7.00 & \PeriodMercury\\
Venus   & \num{4.8675e24} & 0.723 & 0.007 & 3.39 & \PeriodVenus\\
Earth   & \num{5.9722e24} & 1.000 & 0.017 & 0.00 & \PeriodEarth\\
Mars    & \num{6.4171e23} & 1.524 & 0.093 & 1.85 & \PeriodMars\\
Jupiter & \num{1.8982e27} & 5.203 & 0.049 & 1.30 & \PeriodJupiter\\
Saturn  & \num{5.6834e26} & 9.555 & 0.057 & 2.49 & \PeriodSaturn\\
Uranus  & \num{8.6810e25} & 19.22 & 0.046 & 0.77 & \PeriodUranus\\
Neptune & \num{1.02413e26} & 30.11 & 0.011 & 1.77 & \PeriodNeptune\\
\bottomrule
\end{tabular}
\end{table}

\begin{figure}[H]
\centering
\includegraphics[width=\textwidth]{solar_orbits}
\caption{Solar System: inner planets over two years (left) and the whole system over one
orbit of Neptune (right).}
\label{fig:solar}
\end{figure}

A log-log fit of the measured periods against the semi-major axes
(Fig.~\ref{fig:kepler}) gives a slope of \SolarSlope, compared with $3/2$ from Kepler's
third law. The small deviations of the outer planets are real: Jupiter and Saturn perturb
each other and their neighbours, and for Uranus and Neptune a single, eccentric orbit
is too short to average the angular speed. Over 165 years the energy is conserved to
\SolarErrE, the momenta to \SolarErrP\ and \SolarErrL, and $Q=\SolarVirial$.

\begin{figure}[H]
\centering
\includegraphics[width=0.75\textwidth]{solar_kepler}
\caption{Kepler's third law: orbital periods measured in the simulation against the
semi-major axis, with the line $T=a^{3/2}$ (years and AU).}
\label{fig:kepler}
\end{figure}

\subsection{Star clusters and the virial theorem}

Finally we follow 400 equal-mass ``stars'' in $N$-body units with softening
$\varepsilon=0.05$ and $h=0.002$. A \emph{Plummer sphere}, sampled from its equilibrium
distribution function, is born virialized and stays near $Q=1$ (late-time average
\PlummerQlate), fluctuating because $N$ is finite.

A \emph{cold collapse} starts from a homogeneous sphere with every star at rest ($Q=0$).
Gravity makes it collapse in about one free-fall time
$t_{ff}=\frac\pi2\sqrt{R^3/2GM}\approx1.11$; at maximum compression $Q$ overshoots to
\ColdQmax, the cluster rebounds, and within a few $t_{ff}$ it settles into a
quasi-stationary state with $\langle Q\rangle=\ColdQlate$ (Fig.~\ref{fig:cluster}). This
\emph{violent relaxation} (Lynden-Bell, 1967) happens on a dynamical time scale, without
star--star collisions, and is believed to shape elliptical galaxies. The Lagrangian
radii show a dense core forming while the outer 10\,\% of the mass expands into a halo,
part of which escapes (Fig.~\ref{fig:collapse}). The softened energy is conserved to
\ColdErrE\ despite the violent compression.

\begin{figure}[H]
\centering
\includegraphics[width=\textwidth]{cluster_virial}
\caption{Left: virial ratio $Q=2K/|U|$ for the cold collapse and the Plummer sphere.
Right: Lagrangian radii containing 10, 50 and 90\,\% of the mass during the cold collapse.}
\label{fig:cluster}
\end{figure}

\begin{figure}[H]
\centering
\includegraphics[width=\textwidth]{cluster_snapshots}
\caption{Snapshots of the cold collapse ($xy$ projection): initial sphere, contraction,
maximum compression and the final relaxed cluster with escaping stars.}
\label{fig:collapse}
\end{figure}

\subsection{Videos}

The repository also contains 3D videos (MP4) of these experiments, rendered with a
rotating camera and fading trails (\texttt{videos/}). Figure~\ref{fig:stills} shows four
frames.

\IfFileExists{figures/video_stills.png}{
\begin{figure}[H]
\centering
\includegraphics[width=\textwidth]{video_stills.png}
\caption{Frames from the videos: figure-eight choreography, chaos in the six-body system,
the Solar System over one orbit of Neptune and the cold collapse of a star cluster.}
\label{fig:stills}
\end{figure}
}{}

% ===========================================================================
\section{Discussion}
% ===========================================================================

Table~\ref{tab:summary} collects the conservation errors of all runs. Linear and angular
momentum are conserved to machine precision in every case, as expected from a method
built on antisymmetric central pair forces. The energy error is always bounded and its
size is controlled by the ratio of the time step to the shortest dynamical time: it is
largest for the eccentric orbits of the four-body system and for the dense core of the
collapsing cluster.

\begin{table}[H]
\centering
\caption{Maximum conservation errors and virial ratio for every run (velocity Verlet).}
\label{tab:summary}
\begin{tabular}{lcccc}
\toprule
Run & $|\Delta E/E_0|$ & $|\Delta\vect P|$ (rel.) & $|\Delta\vect L|$ (rel.) & $Q$\\
\midrule
Sun--Earth--Moon (2 yr) & \SemErrE & \SemErrP & \SemErrL & \SemVirial\\
Four bodies (15 yr)     & \FourErrE & \FourErrP & \FourErrL & \FourVirial\\
Six bodies (10 yr)      & \ChaosErrE & \ChaosErrP & \ChaosErrL & ---\\
Figure-eight (10 $T$)   & \EightErrE & --- & --- & \EightVirial\\
Solar System (165 yr)   & \SolarErrE & \SolarErrP & \SolarErrL & \SolarVirial\\
\bottomrule
\end{tabular}
\end{table}

The virial theorem is satisfied to $10^{-4}$ or better for all bound, bounded systems. It
does not apply to the six-body system, where bodies escape, nor at early times of the
cold collapse, which is precisely what makes the latter a good illustration of
relaxation towards virial equilibrium.

Compared with the original 2023 work, three corrections are worth noting. First, the
Moon's mass in Table~1 of that report (\SI{2e40}{kg}) was a typo for \SI{7.342e22}{kg}.
Second, the Earth's orbit has eccentricity close to zero, not to one. Third, the orbits of
the four-body system are not circular but have $e=0.5$, as explained by the vis-viva
equation. In addition, measuring errors relative to the initial value (rather than
subtracting the mean) reveals the slow round-off drift of $\vect L$ that the mean hides.

The main limitations are those of any fixed-step direct-summation code. Close encounters
(for example Burrau's Pythagorean three-body problem) require either adaptive time steps,
which must be designed carefully to preserve symplecticity, or regularization. The
$\mathcal O(N^2)$ cost limits the method to a few thousand bodies; tree codes
($\mathcal O(N\log N)$) or fast multipole methods are needed beyond that.

% ===========================================================================
\section{Conclusions}
% ===========================================================================

\begin{itemize}
  \item The velocity Verlet method conserves linear and angular momentum to machine
        precision and keeps the energy error bounded over thousands of orbits, because it
        is symplectic. Non-symplectic methods, even of higher order, drift.
  \item The measured orders of convergence (\OrderVerlet\ for Verlet, \OrderYoshida\ for
        Yoshida-4) agree with theory, and Yoshida-4 offers the best accuracy per cost while
        remaining symplectic.
  \item Analytical results are reproduced quantitatively: Kepler's third law (slope
        \SolarSlope), the eccentricity predicted by the vis-viva equation (\FourEcc), the
        linear instability rate of Lagrange's triangle (\LagSigma\ vs.\ \LagSigmaTh), and
        the virial theorem.
  \item The $N$-body problem is chaotic in general: the six-body system has a Lyapunov
        time of about \ChaosLyapDays\ days, and even the integration step changes the
        trajectory after a few Lyapunov times. Long-term predictions are only meaningful in
        a statistical sense, which is guaranteed by the conservation laws and the
        shadowing property.
\end{itemize}

% ===========================================================================
\appendix
\section{Reproducing the results}
% ===========================================================================

\begin{verbatim}
git clone https://github.com/jesus-garcia-izquierdo/N-body-problem.git
cd N-body-problem
pip install -e .
python examples/01_sun_earth_moon.py      # ... up to 08_cluster_virial.py
python tools/build_documents.py           # regenerates this report
\end{verbatim}

% ===========================================================================
\begin{thebibliography}{99}
\small
\bibitem{aarseth} S.~J. Aarseth, \emph{Gravitational N-Body Simulations}, Cambridge University Press (2003).
\bibitem{benettin} G.~Benettin, L.~Galgani, A.~Giorgilli, J.-M. Strelcyn, \emph{Meccanica} \textbf{15}, 9--30 (1980).
\bibitem{binney} J.~Binney, S.~Tremaine, \emph{Galactic Dynamics}, 2nd ed., Princeton University Press (2008).
\bibitem{chenciner} A.~Chenciner, R.~Montgomery, \emph{Annals of Mathematics} \textbf{152}, 881--901 (2000).
\bibitem{goldstein} H.~Goldstein, C.~Poole, J.~Safko, \emph{Classical Mechanics}, 3rd ed., Addison-Wesley (2002).
\bibitem{hairer} E.~Hairer, C.~Lubich, G.~Wanner, \emph{Geometric Numerical Integration}, 2nd ed., Springer (2006).
\bibitem{haro} G.~Haro Pérez, J.~García Izquierdo, \emph{Problema de los n-cuerpos}, course report (2023).
\bibitem{laskar} J.~Laskar, \emph{Nature} \textbf{338}, 237--238 (1989).
\bibitem{lyndenbell} D.~Lynden-Bell, \emph{Monthly Notices of the RAS} \textbf{136}, 101 (1967).
\bibitem{moore} C.~Moore, \emph{Physical Review Letters} \textbf{70}, 3675 (1993).
\bibitem{szebehely} V.~Szebehely, C.~F. Peters, \emph{Astronomical Journal} \textbf{72}, 876 (1967).
\bibitem{verlet} L.~Verlet, \emph{Physical Review} \textbf{159}, 98 (1967).
\bibitem{yoshida} H.~Yoshida, \emph{Physics Letters A} \textbf{150}, 262--268 (1990).
\end{thebibliography}

\end{document}
